\section{Large Language Models for Code}
\label{sec:bg-llms}

Every generator and judge in this thesis is a transformer~\cite{vaswani2017attention}.
Two properties of the architecture matter here: salience is learned, so a fact
placed in a long prompt competes with everything else, and attention costs grow
with sequence length, making the question of \emph{what} to put in the context
window unavoidable.

Applying transformers to code began with code-specific pre-training.
CodeBERT learns joint representations of natural language and programming
language~\cite{feng2020codebert}; GraphCodeBERT adds a data-flow objective,
and reports gains on downstream code tasks from incorporating that
structure~\cite{guo2020graphcodebert}. General-purpose
instruction-following models such as GPT-4~\cite{openai2023gpt4} subsequently
displaced fine-tuning for many software-engineering
tasks~\cite{hou2023llmse,fan2023llmse}.

Two operational properties bear on the experiments. Decoding is stochastic:
setting temperature to zero requests greedy decoding, which is nearly but not
exactly reproducible. And a model's input is bounded, so selecting a subset of a
repository is a requirement, not an optimisation. The two established ways of
making that selection, similarity-based retrieval and structural traversal,
are the subjects of the next two sections.
